\begin{textsample}{POS Dim 1 – sc – Score 64.00 – sc/sc\_inf\_e\_ef\_29.txt}  \label{ex:f1_pos_015}
Arte 4.2.1 Texto \textbf{introdutório} > O menino era ligado em Despropósitos.
Quis montar os alicerces de uma casa sobre orvalhos.
A mãe reparou que o menino gostava mais do vazio do que do cheio.
Falava que os vazios são maiores e até infinitos [... ]. > ( BARROS, 1999, n.p. ).
A Arte é um artefato da cultura humana e das relações que o sujeito estabelece com o contexto, com os outros sujeitos e com ele mesmo.
A educação em Arte está, pois, ligada à história das culturas da \textbf{humanidade}, que seguiu um padrão \textbf{hegemônico} até as transformações advindas com os pressupostos da modernidade.
Esse movimento trouxe experiências significativas para o ensino da Arte, de modo a ampliar as possibilidades de ensinar e de aprender.
A Arte no currículo da Educação Básica é \textbf{imprescindível}, uma vez que algumas habilidades são próprias dessa área, tais como : a produção artística, a fruição e o conhecimento sensível, que \textbf{agrega} os elementos da percepção, da imaginação, da criação, da \textbf{intuição} e da emoção.
Vale ressaltar que a legitimidade da Arte como \textbf{disciplina} curricular teve um percurso iniciado pela Lei de Diretrizes e Bases da Educação Nacional, Lei Nº 5.692, de 11 de agosto de 1971, com a nomenclatura de Educação Artística, entendida, naquele momento, apenas como “ atividade educativa ”.
Somente com a Lei de Diretrizes e Bases da Educação Nacional Nº 9.394, de 20 de dezembro de 1996, é que a Arte passa a ser considerada obrigatória na Educação Básica, como esclarece o seu Art. 26, §2 : “ O ensino da arte constituirá componente curricular obrigatório, em diversos níveis da educação básica, de forma a promover o desenvolvimento cultural dos \textbf{alunos} ” ( BRASIL, 1996, n.p. ).
É importante destacar que a mesma Lei, no Art. 26, sofreu modificações em 2016, quando esclarece que o ensino e a aprendizagem de Arte nos currículos das escolas passam a incluir artes visuais, dança, música e teatro.
A partir da Lei Nº 9.394/1996 e suas alterações, houve movimentos no sentido de ampliar os cursos de formação docente nessas linguagens, evitando equívocos, como a exigência da polivalência no âmbito escolar, impossibilitando, por vezes, a garantia das manifestações para reconhecimento e legitimidade das linguagens artísticas específicas.
A formação artística e estética do estudante perpassa pela experiência visual, espacial e tátil, pelo movimento corporal, pela expressão corporal no tempo e no espaço e pela manipulação e criação de sons, a partir de um olhar \textbf{crítico}.
Dessa forma, justifica -se a importância dessas linguagens, pois estas, embora tenham suas especificidades, dialogam de forma articulada.
Nessa perspectiva, a relação da cultura e suas diversidades, como as questões étnico-raciais, educação ambiental, \textbf{tecnologia}, entre outras, fortalece a visão integrada de mundo ( FREIRE, 1987 ).
A Base Nacional Comum Curricular ( BNCC ) \textbf{reitera} esse \textbf{posicionamento}, especialmente quando propõe > [... ] que a \textbf{abordagem} das linguagens articule seis dimensões do conhecimento que, de forma indissociável e simultânea, caracterizam a singularidade da experiência artística.
Tais dimensões perpassam os conhecimentos das Artes visuais, da Dança, da Música e do Teatro e as aprendizagens dos \textbf{alunos} em cada contexto social e cultural.
Não se trata de eixos temáticos ou \textbf{categorias}, mas de linhas maleáveis que se interpenetram, constituindo a especificidade da construção do conhecimento em Arte na escola.
Não há nenhuma hierarquia entre essas dimensões, tampouco uma ordem para se trabalhar com cada uma no campo pedagógico. ( BRASIL, 2017, p. 194, \textbf{grifo} do \textbf{autor} ).
As dimensões \textbf{conceituais} do ensino e aprendizagem da Arte na BNCC ( 2017 ) — criação, \textbf{crítica}, estesia, expressão, fruição e reflexão — \textbf{demandam} discussões no âmbito da formação inicial e continuada para que, de fato, sejam tratadas como linhas permeáveis, no sentido de contemplar questões \textbf{críticas}, éticas, estéticas, políticas e culturais, em diálogo com as propostas já existentes em Santa Catarina.
A dimensão da criação concentra -se em uma atitude investigativa, conferindo materialidade estética às ideias inventivas dos estudantes.
Nesse interim, trata das inquietações, dos conflitos e da tomada de decisões nas práticas artísticas, estéticas e culturais, \textbf{tanto} do estudante \textbf{aprendiz}, individualmente ou em seu coletivo, quanto do professor ( BRASIL, 2017 ).
A dimensão \textbf{crítica} na BNCC ( 2017 ) indica a necessidade das relações entre as experiências dos estudantes com as manifestações artísticas e culturais, provocando -lhe um estranhamento do mundo, o que o impulsiona a apropriar -se de novas \textbf{compreensões} do território em que está inserido.
Pela investigação é possível articular a \textbf{crítica} a uma ação e um pensamento propositivo, articulando o ensinar e o aprender em Arte em seus aspectos : políticos, históricos, filosóficos e sociais.
A dimensão estesia trata da experiência sensível dos estudantes no que \textbf{diz} respeito ao espaço/tempo, relacionados ao som, à imagem, ao corpo, bem como as suas materialidades.
Possibilita ao estudante o conhecimento de si mesmo, do outro e do mundo, tornando o corpo protagonista da experiência.
A expressão, manifestada de forma individual e/ou coletiva, atravessa a experiência artístico/estética a partir das linguagens da Arte.
Quando o estudante constrói uma relação de prazer, de estranhamento e de reflexão com o objeto/espaço/obra observado, ativa suas sensibilidades, levando -o a fazer conexões com épocas, pessoas e lugares, afetando -o singularmente e ao seu entorno.
A essa experiência, a BNCC ( 2017 ) chama de “ fruição ”.
A BNCC apresenta 9 \textbf{competências} específicas para o componente Arte, entendida aqui como direitos de aprendizagem.
BASE NACIONAL COMUM CURRICULAR \textbf{COMPETÊNCIAS} ESPECÍFICAS DE ARTE PARA O ENSINO FUNDAMENTAL 1.
Explorar, conhecer, fruir e analisar criticamente práticas e produções artísticas e culturais do seu entorno social, dos povos indígenas, das comunidades tradicionais brasileiras e de diversas sociedades, em distintos tempos e espaços, para reconhecer a arte como um fenômeno cultural, histórico, social e sensível a diferentes contextos e dialogar com as diversidades. 2.
Compreender as relações entre as linguagens da Arte e suas práticas integradas, inclusive aquelas possibilitadas pelo uso das novas \textbf{tecnologias} de informação e comunicação, pelo cinema e pelo audiovisual, nas condições particulares de produção, na prática de cada linguagem e nas suas articulações. 3.
Pesquisar e conhecer distintas \textbf{matrizes} estéticas e culturais – especialmente aquelas manifestas na arte e nas culturas que constituem a identidade brasileira –, sua tradição e manifestações \textbf{contemporâneas}, reelaborando -as nas criações em Arte. 4.
Experienciar a ludicidade, a percepção, a expressividade e a imaginação, ressignificando espaços da escola e de fora dela no âmbito da Arte. 5. \textbf{Mobilizar} recursos tecnológicos como formas de registro, pesquisa e criação artística. 6.
Estabelecer relações entre arte, mídia, \textbf{mercado} e consumo, compreendendo, de forma \textbf{crítica} e problematizadora, modos de produção e de circulação da arte na sociedade. 7.
Problematizar questões políticas, sociais, \textbf{econômicas}, científicas, tecnológicas e culturais, por meio de exercícios, produções, intervenções e apresentações artísticas. 8.
Desenvolver a autonomia, a \textbf{crítica}, a \textbf{autoria} e o trabalho coletivo e colaborativo nas artes. 9.
Analisar e valorizar o patrimônio artístico nacional e internacional, material e imaterial, com suas histórias e diferentes visões de mundo.
No processo de ensino e de aprendizagem em Arte, o professor tem um papel \textbf{relevante}, especialmente no diálogo com os estudantes, oportunizando espaços para que ele tenha a sua \textbf{autoria} valorizada em todos os seus percursos imagéticos e criativos.
Vale destacar que este documento teve a participação de representantes de escolas municipais, estaduais e particulares, resultando na troca de conhecimentos, de experiências e de práticas educativas.
É com essa responsabilidade, de uma educação voltada à formação integral do estudante como ser social, educacional e sensível, que nos empenhamos com essas parcerias a trilhar por percursos de linhas permeáveis, destacando poéticas pessoais, coletivas e, sobretudo, colaborativas.
Fonte : BRASIL.
MEC.
Base Nacional Comum Curricular, 2018. 4.2.2 Metodologia Todo o currículo \textbf{agrega} escolhas epistemológicas e \textbf{metodológicas} coerentes com o que se pensa sobre os sujeitos, seu entorno e as relações construídas.
Dessa forma, o Currículo do Território Catarinense foi tecido levando em conta a Base Nacional Comum Curricular, aprovada e legitimada em dezembro de 2017.
Contudo, vale ressaltar que o Estado de Santa Catarina tem um percurso histórico e cultural nos seus currículos, que necessita da \textbf{compreensão}, do respeito e das (re)significações nos percursos que vão se modificando e alinhando ao longo do tempo.
Destacamos que a metodologia se constitui em uma importante etapa do currículo, o qual está em constante diálogo com os processos de ensino e de aprendizagem em Arte, pois \textbf{sinaliza} alguns \textbf{caminhos} necessários para as práticas educativas.
Neste documento, o componente curricular Arte possui Unidades Temáticas : Artes Visuais, Dança, Música e Teatro, entendendo -as não apenas como temas, mas, sobretudo, como linguagem, expressão e conhecimento.
Além disso, este documento inclui a Unidade Temática “ Artes Integradas ”, que explora as relações e as articulações entre as linguagens, as demais áreas de conhecimento e suas práticas ( BRASIL, 2017 ).
O Quadro 1 a seguir apresenta a Arte no currículo do Ensino Fundamental ( anos iniciais e anos \textbf{finais} ), organizando -a em blocos de habilidades e conteúdos, com ênfase em alguns aspectos.
Quadro 1 - Arte no currículo do Ensino Fundamental PERCURSOS FORMATIVOS EM ARTE | Blocos | Turmas/Anos | Ênfase | |---|---|---| | Bloco 1 | 1º e 2º | Alfabetização em Arte | | Bloco 2 | 3º, 4º e 5º | Arte e cultura local, regional e catarinense | | Bloco 3 | 6º e 7º | Arte e cultura nacional e internacional | | Bloco 4 | 8º e 9º | Arte \textbf{contemporânea} | Fonte : Elaborado pelos \textbf{autores} com base na BNCC ( BRASIL, 2017 ).
A escolha por blocos justifica -se, uma vez que compreendemos que as atividades cognitivas, sensíveis e culturais dos estudantes se constituem e se ampliam \textbf{ancorados} nas experiências e nas culturas, \textbf{permeando} seu cotidiano.
Por conta disso, a ênfase dada para cada um dos blocos fundamentou -se, principalmente, por considerar as trajetórias dos estudantes, respeitar suas singularidades e promover o trabalho colaborativo.
Os \textbf{quadros} apresentados no Apêndice A são compostos pelos seguintes elementos : “ Unidades temáticas ”, “ Objetos de conhecimento ”, “ Habilidades e conteúdos ”.
A Unidade temática é composta por um arranjo dos objetos de conhecimentos no percurso do Ensino Fundamental, adequados às linguagens da Arte.
No que se refere aos objetos do conhecimento em Arte, destacam -se : Contextos e práticas ; Elementos da linguagem ; \textbf{Matrizes} estéticas e culturais ( em Artes Visuais ) ; Materialidades ( em Artes Visuais e Música ) ; Processos de criação ; Sistemas de linguagem ; Notação e registro musical ( em Música ).
As habilidades \textbf{dizem} respeito às aprendizagens essenciais que oportunizem aos estudantes do Ensino Fundamental a formação integral.
Já os conteúdos destacam conceitos e práticas em diálogo com as habilidades e o objeto de conhecimento, que partem das seis dimensões : criação, \textbf{crítica}, estesia, expressão, fruição e reflexão.
Nos \textbf{quadros} do Apêndice A, nas colunas que apresentam os conteúdos, alguns itens encontram -se sublinhados para destaque, em cada bloco, de modo a respeitar o processo do percurso dos estudantes e considerar que as habilidades do componente Arte estão organizadas em anos iniciais e \textbf{finais}.
Entretanto, o professor tem a possibilidade de desenvolver todos os itens ou ampliá -los, visto que estes podem se \textbf{aproximar} mais da sua realidade e dos contextos culturais locais.
No desenvolvimento \textbf{metodológico} do componente Arte, \textbf{sugere} -se o trabalho por projetos, a partir de linhas permeáveis que se conectam entre as linguagens e as Artes Integradas, dando ênfase às culturas locais, regionais e do Estado, nas seguintes modalidades : Educação de Jovens e Adultos, Educação Ambiental, Educação do Campo, Educação Indígena, Educação Quilombola e Educação Especial.
Outra questão a ser considerada sobre o processo \textbf{metodológico} \textbf{diz} respeito também às identidades dos espaços.
Para o componente Arte, é fundamental demarcar no ambiente escolar o seu lugar, o que possibilita ao professor e aos estudantes experiências com suportes, materiais, instrumentos e variados espaços, de forma a nutrir seus processos de criação e de reflexão estética e possibilitar suas produções pessoais, coletivas e colaborativas.
É fundamental, também, um espaço considerável para as \textbf{tecnologias} no processo de ensino e de aprendizagem em Arte, visto que o contexto \textbf{atual} é \textbf{permeado} pelo universo \textbf{digital}, o que possibilita uma outra forma de aprender e de ampliar o conhecimento com olhares e outras percepções.
Embora as linguagens tenham suas especificidades, no contexto \textbf{contemporâneo}, isoladamente elas \textbf{perdem} o sentido plural.
É importante que aconteçam atravessamentos entre as linguagens da Arte, de modo a garantir linhas permeáveis, que perpassam pela pesquisa, pelo conhecimento e por novas descobertas e invenções.
Para a apropriação desta proposta, é necessário pensar -se no percurso do estudante desde a Educação Infantil até o Ensino Médio da Educação Básica, com o \textbf{intuito} de evidenciar um sujeito autônomo, \textbf{crítico}, ético, autoral, sensível, criativo, que trabalhe coletiva e colaborativamente.
Nesse processo, \textbf{espera} -se que o professor construa uma postura ética, estética, investigativa, criativa e, principalmente, que esteja aberto a novas ideias e percepções de ser e de mundo. 4.2.3 Processos de ensinar e aprender na \textbf{disciplina} de Arte A avaliação nos processos de aprendizagem em Arte requer conhecimento, experiência, critérios coerentes e clareza para quem avalia e é avaliado.
Nesse documento, a avaliação é entendida como um processo formativo, com o envolvimento da comunidade escolar, em que o papel do professor é o de ensinar e aprender, mantendo sempre a postura de \textbf{pesquisador}.
Sobre a questão \textbf{avaliativa}, pauta -se na legislação - a Lei Nº 9.394/1996, no Art. 24, inciso V, o qual indica que a avaliação precisa ser “ [... ] contínua e cumulativa do desempenho do \textbf{aluno}, com prevalência dos aspectos \textbf{qualitativos} sobre os quantitativos e dos resultados ao longo do período sobre os de eventuais \textbf{provas} \textbf{finais} ” ( BRASIL,1996, n.p. ).
Nessa perspectiva, a avaliação envolve as questões \textbf{teóricas} e \textbf{metodológicas}, enfatizando a experiência estética, as relações de afeto, a poética pessoal e a produção de conhecimento cooperativa.
Hernández ( 2000 ), na mesma linha de pensamento, compreende a avaliação como um conjunto de ações que \textbf{agrega} uma gama diversificada de fazeres, de instrumentos e de critérios \textbf{avaliativos}.
No entanto, é importante ter -se em \textbf{mente} que a Arte no Ensino Fundamental é a experiência com o sensível de cada humano.
Assim sendo, como afirma Gonçalves ( 2010, p. 164 ), é preciso refletir se “ [... ] aquilo que se está propondo propicia uma aproximação do sujeito ao humano e à ampliação criativa e \textbf{crítica} de suas possibilidades de [ se ] expressar ”.
Entende -se ser de \textbf{suma} importância refletir sobre a avaliação em Arte neste documento, por tratar -se de um processo que \textbf{auxilia} \textbf{tanto} o professor que pesquisa e ensina, quanto o estudante, sujeito \textbf{central} desse processo.
Dessa forma, de posse dessas informações, o professor pode planejar com propriedade suas propostas pedagógicas ; afinal, como afirma Mödinger et al. ( 2012 ), a avaliação é uma das âncoras dos processos de ensinar e de aprender e tem como premissa > [... ] acompanhar, questionar, \textbf{instigar} e principalmente provocar mudanças.
É retomar o que foi feito desde o primeiro momento e gestar novo planejamento com base na observação sistemática e no registro consciente, nos acertos e desacertos, costurando um processo no outro de forma \textbf{dialética}.
Uma avaliação é sempre o embrião da próxima ação pedagógica. ( MÖDINGER et al., 2012, p. 149 ).
Na perspectiva formativa, vale \textbf{salientar} alguns instrumentos de avaliação, a saber : documentos ( anotações pessoais e trabalhos pontuais ), \textbf{seminários}, trabalhos individuais e coletivos, \textbf{exposições}, portfólio, concertos musicais, criação de partituras, performances, protocolos, apresentações teatrais e de dança e exercícios poéticos ( sonoros, corporais e visuais ), autoavaliação, dentre outros.
A autoavaliação, do ponto de \textbf{vista} do estudante, contribui para que este perceba seus percursos de aprendizagem, com o objetivo de identificar avanços e dificuldades a serem suprimidas no diálogo com o professor e demais estudantes.
Para o professor, esse instrumento \textbf{sinaliza} e pode (re)significar suas práticas pedagógicas.
Para Álvarez Méndez ( 2002, p. 87-88 ), “ [... ] quem ensina precisa continuar aprendendo com e sobre sua prática de ensino.
Quem aprende precisa continuar aprendendo constantemente [... ] ”, de modo a perceber as potencialidades dos estudantes e mobilizá -los para buscar e elaborar novos conhecimentos, saberes e experiências em/com Arte.
Outro instrumento que merece destaque é a construção de portfólios, pois eles apresentam evidências dos processos de aprendizagem do estudante e de ensino do professor.
É \textbf{imprescindível} que o estudante identifique “ [... ] o que sabe e o que não sabe, capaz de realizar escolhas, respeitado no seu julgamento que é parte do processo e, mais importante, sendo visto na sua singularidade ” ( OLIVEIRA ; ELLIOT, 2012, p.34 ).
No portfólio, é preciso garantir a \textbf{autoria} e a autonomia do estudante, que pode atravessar e/ou articular as linguagens, suas experiências na escola e fora dela, bem como suas investigações e descobertas.
Assim como os processos de ensinar e aprender estão balizados também nos processos \textbf{avaliativos} com a escolha de seus instrumentos, os critérios têm a mesma relevância.
Nesse sentido, o diálogo entre professor e estudantes, para delinear claramente instrumentos e critérios, fará toda a diferença, \textbf{tanto} nos aspectos cognitivos quanto nos sensíveis.
Para Mödinger et al. ( 2012 ), os critérios \textbf{avaliativos} precisam estar em consonância com aspectos essenciais, os quais devem ser levados em consideração nos processos de aprendizagem em Arte, como : > O comprometimento do [ estudante ] com as discussões e \textbf{tarefas} designadas ; > A participação efetiva em todo o processo que ocorre em sala de \textbf{aula} ; A disponibilidade para pesquisar, investigar e compartilhar conhecimentos e experiências ; A autonomia para expor ideias e \textbf{inter-relacionar} conceitos, conteúdos e produções artísticas ; O cumprimento de prazos estipulados para a entrega ou apresentação de trabalhos ; O respeito mútuo às manifestações dos colegas. > ( MÖDINGER et al., 2012, p. 143 ).
É importante \textbf{lembrar} que esse documento prima por uma avaliação democrática, pois entende -se que o professor de Arte poderá ser um agente mobilizador de ação, de reflexão, de afetos e de diálogos com a vida.

% matched lemmas: abordagem, agregar, aluno, ancorar, aprendiz, aproximar, atual, aula, autor, autoria, auxiliar, avaliativo, caminho, categoria, central, competência, compreensão, conceitual, contemporâneo, crítico, demandar, dialético, digital, disciplina, dizer, econômico, esperar, exposição, final, grifo, hegemônico, humanidade, imprescindível, instigar, inter-relacionar, introdutório, intuito, intuição, lembrar, matriz, mente, mercado, metodológico, mobilizar, perder, permear, pesquisador, posicionamento, prova, quadro, qualitativo, reiterar, relevante, salientar, seminário, sinalizar, sugerir, sumo, tanto, tarefa, tecnologia, teórico, vista
\end{textsample}
